\section{Teorijska podloga}
\label{sec:teorijska-podloga}

\subsection{Segmentacija instanci}
Segmentacija instanci zadatak je računalnog vida koji kombinira detekciju objekata i semantičku segmentaciju. Za razliku od detekcije objekata, koja proizvodi samo okvire (\emph{bounding box}), i semantičke segmentacije, koja svakom pikselu pridružuje oznaku klase bez razlikovanja pojedinačnih objekata, segmentacija instanci za svaku pojedinu instancu objekta u slici proizvodi i okvir i masku na razini piksela. Svakom pikselu unutar maske pridružuje se vrijednost pouzdanosti koja izražava vjerojatnost da pripada detektiranoj instanci.

\subsection{Rubna umjetna inteligencija i kvantizacija modela}
Rubna umjetna inteligencija (\emph{edge AI}) odnosi se na izvođenje modela umjetne inteligencije izravno na ugrađenim ili rubnim uređajima, poput Raspberry Pi računala, umjesto oslanjanja na oblak ili poslužiteljsko računalstvo. Korišteni modeli u pravilu su istovjetni svojim punoopsežnim inačicama, uz ključnu razliku da su kvantizirani. Kvantizacija je postupak smanjenja preciznosti težina modela, tipično s 32-bitnog zapisa s pomičnim zarezom na 8-bitni cjelobrojni zapis, čime se znatno smanjuje potrošnja memorije i računalno opterećenje, omogućujući zaključivanje u stvarnom vremenu na maloj snazi.

\subsection{Optički tok i praćenje značajki}
Optički tok temelji se na praćenju istih značajki kroz više uzastopnih slika snimljenih u različitim trenucima. Pretpostavlja se da se osvjetljenje značajke ne mijenja bitno između slika. Značajke koje se prate prvo se identificiraju detektorom kutova poput FAST ili Harris, koji pronalazi točke u slici dovoljno izražene da se mogu pouzdano pratiti. Lokalni prozor oko svake značajke potom se koristi za pronalaženje njezinog odgovarajućeg položaja u sljedećoj slici minimizacijom razlike u intenzitetu piksela. Ovakav pristup sužava prostor pretrage i čini postupak izvedivim u stvarnom vremenu.

\subsection{Vizualno serviranje temeljeno na slici}
Vizualno serviranje temeljeno na slici (IBVS) metoda je upravljanja koja slikovnu informaciju izravno koristi kao ulaz u upravljačku petlju. Pogreška se definira kao udaljenost u pikselima između detektirane točke prijanjanja i središta slike. Ta se pogreška dovodi proporcionalnom regulatoru koji generira naredbu brzine za letjelicu, vodeći točku prijanjanja prema središtu slike. Upravljačka petlja izvodi se kontinuirano: svaka iteracija preuzima novu sliku, ponovno izračunava pogrešku i izdaje novu naredbu brzine sve dok se pogreška ne minimizira.

\section{Sklopovlje}
\label{sec:sklopovlje}

\subsection{Platforma letjelice}
Letjelica korištena u ovom radu izgrađena je na 7-inčnom karbonskom okviru, odabranom zbog krutosti i male mase. Pogon osiguravaju četiri \emph{brushless} motora T-Motor F90 2810 predviđena za propelere veličine 7 do 8 inča. Upravljanje letom obavlja SpeedyBee F405 upravljačka jedinica koja koristi MAVLink protokol za komunikaciju s računalom na letjelici.

U odnosu na konfiguraciju iz diplomskog rada, natjecateljski scenarij pretpostavlja drukčije postavljanje kamere -- kamera je fizički okrenuta prema gore umjesto prema naprijed -- što zahtijeva novi nosač kamere te izmjenu geometrije preslikavanja pogreške iz slikovne ravnine u naredbe letjelici (vidi Odjeljak~\ref{subsec:preslikavanje-osi}).

\subsection{Računalo i senzori na letjelici}
Računalo na letjelici je Raspberry Pi 5 s 8 GB radne memorije, uz Raspberry Pi AI Hat+ dodatak za sklopovski ubrzano zaključivanje modela. Kao izvor slike koristi se Arduino Nicla Vision kamera s integriranim ToF senzorom udaljenosti. Za razliku od diplomskog rada, u ovom se radu ToF očitanje ne koristi za okidanje manevra prijanjanja -- ono je predviđeno za drugu namjenu izvan opsega ovog rada -- pa manevar prijanjanja mora biti okinut na temelju drugog signala (vidi Odjeljak~\ref{subsec:prijanjanje}).

Za prijanjanje se koristi prilagođeni 3D-tiskani pasivni hvataljka s oprugom, izrađena od PETG materijala, montirana na okvir letjelice. Hvataljka ostaje otvorena dok se na unutarnji dio čeljusti ne primijeni dovoljna sila, nakon čega se hvataljka zatvara.

\section{Implementacija sustava}
\label{sec:implementacija}

\subsection{Cjevovod za detekciju}
\label{subsec:cjevovod}

Cjevovod za detekciju grane preuzet je gotovo bez izmjena iz diplomskog rada, budući da je riječ o obradi isključivo u slikovnoj ravnini, neovisnoj o fizičkoj orijentaciji kamere u prostoru. Cjevovod je dizajniran modularno: svi razredi izvora slike dijele zajednički roditeljski razred, a svi koraci naknadne obrade nasljeđuju zajednički razred \texttt{Stage}. Podaci proizvedeni u svakom koraku pohranjuju se u strukturi \texttt{PipelineContext}, čime koraci nisu međusobno čvrsto povezani i mogu se zamjenjivati ili uklanjati bez utjecaja na ostatak cjevovoda.

Praćenje objekata provodi se algoritmom ByteTrack, koji kombinira Kalmanov filtar s Mađarskim algoritmom za povezivanje detekcija kroz uzastopne slike na temelju preklapanja okvira. Nakon toga slijede četiri koraka naknadne obrade: izdvajanje binarne maske pragiranjem izlazne pouzdanosti segmentacije, izračun mape udaljenosti (\emph{distance heatmap}) Euklidskom transformacijom udaljenosti radi određivanja najdebljeg i najstabilnijeg dijela grane, računanje medijalne osi (kostura) binarne maske radi ograničavanja kandidata na središnjicu grane, te konačno bodovanje kandidata težinskom funkcijom koja kombinira pouzdanost detekcije, udaljenost od središta slike i veličinu detekcije.

Nakon početnog razdoblja zagrijavanja (\emph{warmup}) od približno 20 sekundi, tijekom kojeg se prikupljaju najbolji kandidati iz svake slike, prikupljene točke grupiraju se algoritmom hijerarhijskog grupiranja s pragom udaljenosti od 30 piksela. Najpopunjenija skupina odabire se kao konačna točka prijanjanja (\emph{final point}), nakon čega se model detekcije prestaje izvoditi, a praćenje preuzima KLT algoritam opisan u nastavku.

\subsection{Vizualno serviranje temeljeno na slici (IBVS)}
\label{subsec:ibvs}

Nakon zaključavanja konačne točke, IBVS izdvaja FAST i Harris kutne značajke iz referentne slike u lokalnoj okolini točke prijanjanja te pohranjuje pomak svake značajke u odnosu na tu točku. U svakoj sljedećoj slici značajke se prate KLT (\emph{Kanade--Lucas--Tomasi}) algoritmom optičkog toka, a položaj točke prijanjanja procjenjuje se primjenom pohranjenih pomaka na trenutne položaje praćenih značajki. Time se izbjegava ponovno pokretanje modela detekcije za svaku sliku, čime se smanjuje računalno opterećenje i omogućuje praćenje u stvarnom vremenu.

Pogreška se računa kao udaljenost u pikselima između procijenjenog položaja točke prijanjanja i središta slike, zasebno za vodoravnu i okomitu os. Ova komponenta sustava također se preuzima bez izmjena iz diplomskog rada, jer izlaz IBVS-a ostaje par vrijednosti pogreške u slikovnoj ravnini (\texttt{error\_x}, \texttt{error\_y}) bez ugrađene pretpostavke o fizičkoj orijentaciji kamere -- preslikavanje tih vrijednosti u naredbe letjelici u potpunosti je izmješteno u upravljački sloj opisan u nastavku.

\subsection{Zamjena upravljačkog sloja}
\label{subsec:zamjena-upravljanja}

U diplomskom radu upravljanje letjelicom izvodio je namjenski razvijeni ROS čvor koji je koristio OptiTrack odometriju za apsolutno pozicioniranje te trokoračni konačni automat (ALIGNING -- APPROACHING -- PERCHING) u kojem su ispravci iz IBVS-a preslikavani u naredbe apsolutnog položaja: vodoravna pogreška slike preslikavala se u y-os letjelice, okomita pogreška u z-os, dok se x-os (smjer prilaska) pomicala unaprijed fiksnim korakom sve dok ToF senzor ne bi zabilježio udaljenost ispod praga, nakon čega bi se izvela dvokoračna trajektorija prijanjanja.

Budući da natjecateljski scenarij ne raspolaže sustavom OptiTrack niti koristi ToF senzor za okidanje prijanjanja, taj se upravljački čvor u potpunosti zamjenjuje otvorenim paketom \texttt{ibvs\_perching}~\cite{ibvs_perching}. Taj paket ne koristi kontroler pozicije niti vanjsku lokalizaciju: letjelicom upravlja izravno preko sučelja \texttt{mavros/setpoint\_raw/attitude}, šaljući naredbe kutne brzine tijela (\emph{body rate}) za bočno poravnanje te naredbu brzine penjanja (kroz polje \texttt{thrust}, protumačeno u ArduPilot GUIDED načinu rada kao naredba brzine penjanja, a ne kao sirovi potisak) za okomiti prilazak.

Sučelje između percepcije i ovog upravljačkog sloja jest poruka \texttt{PointStamped} na temi \texttt{ibvs/target\_point}, u kojoj \texttt{x} i \texttt{y} predstavljaju normaliziran pomak ciljne točke od središta slike, a \texttt{z} (opcionalno) procijenjenu udaljenost do cilja. Ovime se zadržava jasna granica između percepcijskog cjevovoda (koji ne zna ništa o geometriji letjelice) i upravljačkog sloja (koji ne zna ništa o tome kako je točka detektirana) -- ista podjela odgovornosti koja je postojala i u izvornom sustavu iz diplomskog rada.

\subsubsection{Preslikavanje osi za kameru okrenutu prema gore}
\label{subsec:preslikavanje-osi}

% TODO: ovo poglavlje dovrsiti nakon eksperimentalne provjere stvarnog predznaka i
% osi -- trenutno opisuje namjeravano ponasanje, ne izmjereno

Paket \texttt{ibvs\_perching} izvorno pretpostavlja kameru okrenutu prema dolje, gdje vodoravni pomak u slici (\texttt{point.x}) odgovara pomaku unaprijed u koordinatnom sustavu letjelice, a okomiti pomak u slici (\texttt{point.y}) odgovara pomaku ulijevo/udesno. Budući da je u ovom radu kamera okrenuta prema gore, a letjelica polazi već pozicionirana ispod grane, preslikavanje osi mora se iznova izvesti: os prilaska više nije vodoravna os letjelice (naprijed), nego okomita os (penjanje), dok se poravnanje u preostale dvije osi izvodi analogno prethodnom sustavu.

\subsection{Manevar prijanjanja bez ToF okidača}
\label{subsec:prijanjanje}

% TODO: opisati konacan mehanizam okidanja nakon eksperimentalne provjere

Budući da senzor udaljenosti nije dostupan za okidanje manevra prijanjanja, potreban je drukčiji signal za prijelaz iz faze poravnanja u fazu izvođenja manevra. Kako letjelica polazi s poznate, fiksne udaljenosti ispod grane, razmatra se okidanje na temelju vremena provedenog u stanju poravnanja unutar zadanog praga pogreške, umjesto na temelju izravnog mjerenja udaljenosti.

\subsection{Treniranje modela}
\label{subsec:treniranje-modela}

Model za segmentaciju grana (YOLOv8m-seg) preuzet je bez izmjena iz diplomskog rada. Skup podataka za treniranje sastojao se od fotografija grana snimljenih iz bočne perspektive, a maske za segmentaciju automatski su generirane modelom SAM3 na temelju tekstualnog upita. Model je izvezen u ONNX format i zatim kompiliran u Hailo izvršni format (\texttt{.hef}) za izvođenje na Raspberry Pi AI Hat+ akceleratoru.

% TODO: napomenuti u Raspravi da skup podataka trenutno ne sadrzi fotografije grana
% snimljene iz perspektive kamere okrenute prema gore (nebo kao pozadina), sto moze
% utjecati na tocnost detekcije u novoj konfiguraciji
